\scanpage{110}
\margin{\diagram{\begin{tikzcd}[row sep=small,column sep=small]
Q'(i) \arrow[r,dash] \arrow[d,dash] & L \arrow[d,dash] \arrow[dr,dash] & \deleted{\unsure{Q'(i)}} \\
Q' \arrow[r,dash] & K \arrow[d,dash] \arrow[dr,dash] & L_0 \arrow[d,dash] \\
& Q & K_0
\end{tikzcd}}\unsure{galoisienne}}

On veut construire le groupe de t.m. $G(q)$ associé à l'ens $M(q)$ des $\lambda\in L$ tels que $\lambda\bar\lambda=q^i$, et $\lambda$ unités exc. $p$. On a vu que
\[
M_0(q)/\text{Torsion}\subset (\mathbb Z^P)' \qquad (1)
\]
\margin{N.B. $M_0(q)$ \inserted{$=M_0(p)$} \unsure{gr. des} unités $u\in L$ \unsure{qui sont} exc. $p$ et \unsure{satisf.} $u\bar u=1$.}
$P$ l'ens des idéaux premiers de $L$ sur $p$.

On peut supposer que les idéaux $\mathfrak p_0$ de $K$ se \uline{décomposent} dans $L$. Choisissons un des \deleted{\unsure{ces}} idéaux \unsure{au-dessus de $p$}, $\mathfrak p$, ce qui donne une forme linéaire \unsure{can.} sur $M$
\[
\lambda\longmapsto v_{\mathfrak p}(\lambda),
\]
qui engendre d'ailleurs $\check M$ mod. torsion sous le groupe de Galois $\Gamma$ de $L$, d'où
\diagram{\begin{tikzcd}
G_{mL}\arrow[r,"\nu_L"] & G_L
\end{tikzcd}}

Soit $\Gamma_{\mathfrak p}$ le groupe de décomposition pour $\mathfrak p$, \deleted{\unsure{qui est}} \deleted{\unsure{alors}} et soit \deleted{$K_0$ le corps} $L_0$ le corps des invariants de $\Gamma_{\mathfrak p}$, $K_0=L_0\cap K$. Alors $L_0$ est le plus petit corps de rationalité de $\nu$, qui s'écrit \deleted{donc}
\diagram{\begin{tikzcd}
\nu:G_{mL_0}\arrow[r] & G_{L_0}
\end{tikzcd}}

D'autre part, si $\mathfrak p_0$ est induit par $\mathfrak p$ sur $L_0$, alors \deleted{\unsure{on a l'anneau de valuation associé}} \inserted{l'anneau local complété} de $\mathfrak p_0$ est \deleted{\unsure{égal à celui de}} isom. à celui de $\mathbb Z_{\unsure{\ell|p}}$. Notons que le couple $(L_0,\mathfrak p_0)$ est défini à isom unique près par la donnée de la \uline{classe} de $L$, et peut se définir \deleted{\unsure{comme}} \inserted{via} le \unsure{hensélisé} relatif de $\mathbb Z_{(\unsure{p\mathbb Z|p})}$ dans $L$ \dots. Notons que $L_0$ est purement imaginaire, grâce à l'hypothèse que $\tau\mathfrak p\ne\mathfrak p$.
\mspage{109}
